\documentclass[aspectratio=169]{beamer}

\usetheme{Madrid}
\usecolortheme{seahorse}
\usepackage{graphicx}
\usepackage{booktabs}
\usepackage{siunitx}

\title{AM CVn: Testing a Superhump Interpretation}
\author{TODO: Team Members}
\date{PHYSICS 100 Final Project}

\begin{document}

\begin{frame}
  \titlepage
\end{frame}

\begin{frame}{Scientific Goal}
  \begin{itemize}
    \item AM CVn is an ultracompact helium-accreting binary.
    \item Literature periods: orbital \(1028.7322\) s, positive superhump \(1051.2\) s, strong harmonic near \(525.6\) s.
    \item Period sources: Skillman et al. 1999; Nelemans et al. 2001; Roelofs et al. 2006.
    \item Question: can our 2026 May 20 photometry tighten those constraints?
    \item Rule: separate free-period constraints from literature-conditioned fits.
  \end{itemize}
\end{frame}

\begin{frame}{Observations}
  \begin{columns}
    \begin{column}{0.50\textwidth}
      \begin{itemize}
        \item 140 V-band frames.
        \item 35, 40, and 45 s exposures.
        \item 138 good points after rejecting two artifact frames.
        \item Baseline: 111.20 min.
        \item Covers 6.35 cycles of a 1051.2 s superhump.
      \end{itemize}
    \end{column}
    \begin{column}{0.50\textwidth}
      \includegraphics[width=\textwidth]{analysis/am_cvn/plots/AM_CVn_sequence_X42421BZ_light_curve.png}
    \end{column}
  \end{columns}
\end{frame}

\begin{frame}{Reduction}
  \begin{itemize}
    \item Bias, dark, and star-masked flat-field calibration.
    \item AAVSO sequence X42421BZ used for target and comparison-star geometry.
    \item Fluxes converted to rates before combining different exposure times.
    \item Comparison stars normalized by their own median before forming the ensemble.
    \item Robust target scatter: 1.51 percent; median formal error: 0.90 percent.
  \end{itemize}
\end{frame}

\begin{frame}{Comparison-Star Controls}
  \begin{columns}
    \begin{column}{0.45\textwidth}
      \begin{itemize}
        \item Each comparison star was treated as a pseudo-target.
        \item The pseudo-target was divided by the remaining comparison ensemble.
        \item Robust scatter: 0.70--1.22 percent.
        \item Empirical constant-model \(\chi^2_\nu\) values are near unity.
        \item This supports the ensemble as a stable photometric reference.
      \end{itemize}
    \end{column}
    \begin{column}{0.55\textwidth}
      \includegraphics[width=\textwidth]{analysis/am_cvn/plots/AM_CVn_comparison_star_controls.png}
    \end{column}
  \end{columns}
\end{frame}

\begin{frame}{Models and BIC}
  \scriptsize
  \begin{columns}
    \begin{column}{0.46\textwidth}
      \begin{itemize}
        \item All models fit sine/cosine curves to the same light curve.
        \item Harmonic superhump: fix \(P_{\rm sh}=1051.2\) s and add the \(2f\) harmonic.
        \item Free double wave: same two-frequency shape, but \(P\) is fit from our data.
        \item Beat disk: subtract frequencies, \(f_{\rm sh}=f_{\rm orb}-f_{\rm prec}\), then fit that waveform.
        \item \(\chi^2\): weighted residual size.
        \item \({\rm BIC}_{\rm eff}=\chi^2+k_{\rm eff}\ln N\): fit quality plus complexity penalty.
      \end{itemize}
    \end{column}
    \begin{column}{0.54\textwidth}
      \begin{table}
        \tiny
        \begin{tabular}{lrrr}
          \toprule
          Model & Period & \(k_{\rm eff}\) & \(\Delta{\rm BIC}_{\rm eff}\) \\
          \midrule
          Constant flux & -- & 1 & 96.75 \\
          Single sine & 525.06 s & 4 & 37.44 \\
          1051.2 s + harmonic & fixed & 5 & -4.37 \\
          Beat-constrained disk & 1055.22 s & 6 & 0.00 \\
          Free double wave & 1055.21 s & 6 & 0.00 \\
          \bottomrule
        \end{tabular}
      \end{table}
    \end{column}
  \end{columns}
\end{frame}

\begin{frame}{Model Fits}
  \centering
  \includegraphics[height=0.80\textheight]{analysis/am_cvn/plots/AM_CVn_stacked_model_fits.png}
\end{frame}

\begin{frame}{Constraints}
  \small
  \begin{columns}
    \begin{column}{0.47\textwidth}
      \begin{itemize}
        \item Free single sine: \(525.06\) s; Monte Carlo median \(524.38^{+5.75}_{-3.84}\) s.
        \item Free double wave: \(1055.21^{+13.10}_{-9.63}\) s.
        \item Best independent period measurement: free double wave.
        \item Best conditional model: fixed 1051.2 s + harmonic.
        \item Both support the known superhump family.
      \end{itemize}
    \end{column}
    \begin{column}{0.53\textwidth}
      \includegraphics[width=\textwidth]{analysis/am_cvn/plots/AM_CVn_physical_literature_constraints.png}
    \end{column}
  \end{columns}
\end{frame}

\begin{frame}{Validation}
  \small
  \begin{columns}
    \begin{column}{0.50\textwidth}
      \begin{itemize}
        \item Multi-harmonic Lomb-Scargle recovers the double-wave fundamental.
        \item Fixed-time Monte Carlo perturbations replace row-resampling bootstrap as the primary error analysis.
        \item Beat-constrained model implies \(P_{\rm prec}\approx 11.39\) hr, inferred from the beat relation rather than directly observed.
        \item Leave-one-comparison-star-out period span: 5.2 s.
        \item Outlier clipping period shift: 0.02 s.
      \end{itemize}
    \end{column}
    \begin{column}{0.50\textwidth}
      \begin{itemize}
        \item Even the best sine/double-wave models have reduced chi-squared around 2.5, so uncertainties are empirically scaled and conclusions are conservative.
        \item Residual periodogram has no significant remaining peak.
        \item Systematics terms lower RMS to 1.41\%, but worsen BIC/CV.
        \item Robust losses shift the answer to 1057--1059 s and worsen RMS.
        \item First/second halves differ by 22.8 s.
      \end{itemize}
    \end{column}
  \end{columns}
\end{frame}

\begin{frame}{Could It Be Something Else?}
  \begin{itemize}
    \item Broadband photometry alone cannot prove the object class.
    \item Literature spectroscopy shows helium accretion signatures and an S-wave orbital feature.
    \item A non-binary pulsator would need a spectrum without accretion features and a different multi-period pattern.
    \item The decisive test is time-resolved spectroscopy, not another photometric periodogram.
  \end{itemize}
\end{frame}

\begin{frame}{Conclusion}
  \begin{itemize}
    \item We detect AM CVn short-timescale variability.
    \item The 525 s single-sine signal is best interpreted as a harmonic.
    \item The fixed 1051.2 s harmonic template has the best effective BIC, but it is a literature-conditioned model.
    \item The free double wave is the best independent period measurement.
    \item Detrending, robust likelihoods, and comparison subsets do not improve the constraint.
    \item We do not tighten the literature periods with this dataset.
    \item A longer multi-night baseline is the path to a stronger constraint.
  \end{itemize}
\end{frame}

\end{document}
